\section{Palamidessi reads Evola}

There is an Italian commentary on, and reprinted text of, a letter from Evola to Tomasso Palamidessi here\footnote{\url{http://www.fondazionejuliusevola.it/Documenti/Evola_Palamidessi.pdf}}. Although this is not a translation, this is a rough survey of the contents of the letter, with a transliteration attempted at certain points, to convey or emphasize the meaning.

The commentator notes that the interactions between various schools, personalities in Italian esotericism was \emph{non sempre pacifici} (not always peaceful). Although Palamidessi's career has not been as well studied as Evola's, for obvious metaphysical reasons of ``heft'', this letter sheds a great deal of light. Palamidessi was born in Pisa in 1915, and moved to Turin at the age of 18, in 1933, having already demonstrated an interest in botany, astronomy, astrology, and medicine. After getting to Turin, he dedicated himself \emph{ad formazione occultista integrale} (to the formation of an integrated occultism). One of his teachers was the director of Museo Egizio, Ernest Scamuzzi. Around 1940, he had contact with Alfred Witte\footnote{\url{http://astrologer.ru/Witte/biography_eng.html}}`s Amburgo, as well as with Alexandre Volguine\footnote{\url{http://www.solsticepoint.com/astrologersmemorial/volguine.html}} (French astrologer), and some other personages connected with Masonry or Martinism -- Umberto Gorel Porciatti and Gino Witnesses. With the exception notably of astrology, all of his interests coincide with those of Evola's, and in fact, Palamidessi specifically replies to Evola's Man as Power (1926) \& Hermetic Tradition (1931)\footnote{\url{http://en.wikipedia.org/wiki/The_Hermetic_Tradition}} in pamphlets like The Erotic Power of Kundalini-Yoga (1949). Palamidessi's pamphlets, unlike Evola's metaphysically hefty tomes, have a technical, practical ``cut''. His debt to Evola is quite clear, when focusing on the content, and in some cases, Evola noticed a very real plagiarism. Palamidessi's doctrine follows closely three cardinal points of Evola:
\begin{enumerate}
\item Hermetism \& Alchemy are interpreted in light of tantric yoga

\item There is a hidden initiation thereby available to those who are ``not profane''

\item Tantric yoga is this path
\end{enumerate}

Palamidessi eschews a merely Theosophical orientation towards the ``subtle bodies'' (the lunar path).
\emph{Questa, quindi, la linea di pensiero comune} (This, therefore or thus far, is the line of common thought. The analyst points out that Evola and Palamidessi engaged in ripostes over time, and that Evola's final veredictum was formulated later:

\begin{quotex}
\textit{La donna assoluta non solo non possiede quell'Io, ma non saprebbe nemmeno che farsene, essa non sa nemmeno concepirlo e la sua presenza agirebbe in modo estremamente disturbatore presso a ogni genuina estrinsecazione della di lei più profonda natura.}
\flright{\textsc{J. Evola}, \emph{Metafisica del sesso}, Roma: Mediterranee 1958 , 1996, p. 179}

The absolute woman not only does not possess the ``I'', but she wouldn't even know what to make of it, she does not know how to conceive it, and its presence would act in an extremely disruptive manner in every genuine objectification of her deepest nature.
\end{quotex}


Palamidessi seems to have oscillated back and forth between ``explicit admiration'' for the author of such works he had borrowed heavily from, \& ardent criticism on this one point. Here is the untranslated Italian:


\begin{quotationx}
\textit{Siamo decisamente contro a tutti quei filosofi ed occultisti tanto ideologicamente nemici della donna da relegarla per molti secoli nel numero delle cose e degli esseri inferiori, laddove invece avrebbe potuto essere esaltata quale vivente espressione dell'Eterno femminino. Mi stupisce come un grande scrittore di esoterismo orientale della statura di Julius Evola ... abbia commesso il grave errore di negare alla donna il privilegio dell'immortalità e di quel principio che fa dell'uomo un essere superiore e imperituro. C'è proprio da chiedersi con quale criterio l'Evola abbia definito in un suo articolo sulla rivista Ignis (gen.-feb. 1925, n° 1- 2) ... ``la donna come cosa'' ... . L'uomo e la donna, di qualsiasi sviluppo morale e mentaleessi siano, non sono cose, ma individui, il cui livello evolutivo è tale da manifestare, o no, quel principio che fa dell'individuo un Dio. Ad ogni modo, questa mia constatazione ... non infirma affatto la mia simpatia di studioso nei confronti del dr. Evola che, se ha commesso degli errori di valutazione psicologica nei confronti della donna, ha pure dimostrato una profondità di vedute e di pensiero in altri settori della scienza occulta.}

We decisively oppose all those philosophers and occultists who are ideologically enemies of woman and relegate her for many centuries to be among things and inferior beings, when instead, she should be exalted as the living expression of the Eternal Feminine. It amazes me how a great writer of Eastern metaphysics of the stature of Julius Evola had created the serious error of denying to woman the privilege of immortality and of that quality that makes of man a higher and eternal being. It is proper to ask ourselves by what criteria does Evola define woman as a thing.

Man and woman, of whatever moral and intellectual development, are not things, but individuals, whose evolutionary level is such to manifest, or not, that principle that makes man a god. In every way, my thinking does not invalidate at all my scholarly sympathy toward Dr. Evola who, if he committed some errors of the psychological evaluation toward women, has only demonstrated a depth of insight and thought in other areas of occult science.
\end{quotationx}


It is only fair to note that Palamidessi had engaged in an extensive reading of Evola for years, up to and including virtual ``plagiarism'', as Evola himself at one point noticed (Guenon would later on comment, on the basis of reading excerpts that he was a ``charlatan'', although Evola seems to have wound up on equitable terms with Tomasso, both of whom died about the same time). Evola granted women the \emph{corpi sottili} (subtle body) but not the transcendental I, a position with roots in Weininger.

Here I can't help wondering if there was simply a blundering error on one or both of Evola or Palamidessi's part, over the wisdom application of a principle. The fact that great import or danger falls or turns on such a juncture should not amaze us. Here is an apocryphal\footnote{\url{http://creation.com/battle-quote-not-luther}} Luther quote --

\begin{quotex}
If I profess, with the loudest voice and the clearest exposition, every portion of the truth of God except precisely that little point which the world and the devil are at that moment attacking, I am not confessing Christ, however boldly I may be professing Christianity. Where the battle rages the loyalty of the soldier is proved; and to be steady on all the battle-field besides is mere flight and disgrace to him if he flinches at that one point.
\end{quotex}


And here we can intuit that although there is incredible danger of the ``Eternal Feminine\footnote{\url{http://www.bu.edu/arion/files/2010/03/Paglia-Great-Mother1.pdf}}'' being used in the Dark Age of the Kali Yuga to exalt chaos(as the virtuoso Paglia flirts with doing here) , that perhaps at the higher levels of debate, a very tragic misunderstanding can occur between those of deep insight. This may or may not be the case with Palamidessi.

In fact, it is an old Northern custom\footnote{\url{http://www.mun.ca/mst/heroicage/issues/9/forum2.html}} that a woman can act in the place of a fallen or absent male champion. Tolkien knew\footnote{\url{https://www.youtube.com/watch?v=MSNPeJAgBzo}} about this, and it puts one in mind of Deborah the judge, Esther the queen, and Judith the champion\footnote{\url{http://en.wikipedia.org/wiki/Book_of_Judith}}. The female is all the more potent for being forced to do violence to her sex and transcend it (in this case, taking the male place).

One is also tempted to wonder -- if Satan counterfeits ``traditional'' metaphysics by causing the resurgence of suppressed elements during the Kali Yuga, to the point of them achieving political power, wouldn't that mean he is counterfeiting something real, even in the feminine element?

It may be that Palamidessi, and Evola, as appeared to occur late in real life, can be reconciled.

\emph{More to come next week, with more transliteration from the letter, and with deference \& thanks to Cologero for suggestions.}


\flrightit{Posted on 2012-04-28 by Logres}

\begin{center}* * *\end{center}

\begin{footnotesize}
\begin{sffamily}
\texttt{Cologero on 2012-04-30 at 23:31 said:}

What can you conclude from the notion that a ``woman can act in the place of a fallen or absent male champion.'' The simply reinforces her secondary nature, doesn't it? How would we react to a man acting the the place of a fallen female? Wouldn't we consider that a lessening of his powers rather than an increase?

Why is the ``male gaze'' worthy of capital punishment? Is it simply because Thryth objects to it as a woman, or because she objects to it as queen?

You could have chosen a real example, rather than literary creations, e.g., St Joan of Arc.

I don't understand this: ``if Satan counterfeits ``traditional'' metaphysics by causing the resurgence of suppressed elements during the Kali Yuga, to the point of them achieving political power, wouldn't that mean he is counterfeiting something real, even in the feminine element?''

What elements are suppressed in the Kali Yuga? Doesn't a counterfeit always counterfeit something real? No one creates a knock off of a fake Rolex.

\hfill

\texttt{Cologero on 2012-04-30 at 23:43 said:}

Readers interested in the question about the ``I'' should read Chapter VIII in Weininger's \textit{Sex and Character}.

As for the link to the postmodern journal on Northern customs: there is no disagreement about a continuum of male and female characteristics. However, the point they miss, is that to define the continuum requires two significant pieces of knowledge. To ignore these, is tantamount to expecting a colour-blind man to define the visible light spectrum.

\begin{enumerate}


\item One must know what are the male and female characteristics.

\item One must know the two poles of the spectrum: absolute man and absolute woman.
\end{enumerate}


To try to define these empirically is to beg the question. The only solution is \emph{a priori}.

\hfill

\texttt{Aeneas on 2012-05-01 at 23:41 said:}

If Evola exaggerates the differences between man and women, it seems that Palamidessi wants to obliterate every distinction. What, then, is the male equivalent of the Divine Feminine, if not the Holy Spirit? Or put another way, the male is Purusha and the female Prakriti. But Purusha is Atman (the ``I''), so where does that leave Prakriti? Has Evola betrayed oriental metaphysics, or is he drawing a logical conclusion?

\hfill

\texttt{Charlotte on 2012-05-02 at 06:57 said:}

Holy Soul, the Shekinah, Bride of Christ

\hfill

\texttt{logres on 2012-05-04 at 06:32 said:}

Agreed concerning Joan of Arc. Journal chosen due to Tolkien's scene with the woman -- rather than being ``feminist'' or anti-tradition, it is a reaffirmation of male characteristics. What greater insult than to be beaten by, not the king, but his little daughter? However, there is something ``more'' in Christianity than Evola gives in relation to sex; Evola never discusses to what extent the upsurge of chaos \& the female are actually in contingency with the divine plan; our alternatives are not Evolutionary view vs. Evola. I realize most people here know that, so it was a cheap point. I should concentrate on discerning more along the lines of Charlotte's comment, which was where I am headed. To that end,

\url{http://www.kober.com/marriage.php}

Bo-Yin-Ra

Also, Mouravieff's comments about the knight and his true lady who advance step by step up the stairs of the divine Solfeggio notes in Gnosis obviously is speaking of ascetic practices within and outside marriage designed to realize the Eternal Union of the separated poles. I'll try to excerpt this more directly very soon.

\hfill

\texttt{logres on 2012-05-04 at 06:34 said:}

Palamidessi has a pamphlet directly concerning occult characteristics of the male \& female, so I won't comment anymore on him till after I have obtained \& read.

\hfill

\texttt{Charlotte on 2012-05-04 at 20:13 said:}

essay here you may or may not find interesting, talking about Guenon's astrology, end of Kali Yuga, shift from polar to solar etc

\url{http://www.alignment2012.com/polar-to-solar.html}

Cx

\hfill

\end{sffamily}
\end{footnotesize}

